\documentclass[10pt,a4paper]{article}
\usepackage{fontspec}
\setmainfont{texgyretermes}[Extension=.otf,UprightFont=*-regular,BoldFont=*-bold,ItalicFont=*-italic,BoldItalicFont=*-bolditalic]
\usepackage[margin=1.8cm,headheight=14pt]{geometry}
\usepackage{amsmath,amssymb,booktabs,tabularx,array,ragged2e,enumitem,xcolor,fancyhdr,titlesec,microtype}
\usepackage[hidelinks,unicode]{hyperref}
\definecolor{ink}{HTML}{183A45}\definecolor{teal}{HTML}{087F7A}\definecolor{gray}{HTML}{56656B}
\newcolumntype{P}[1]{>{\RaggedRight\arraybackslash}p{#1}}
\renewcommand{\tabularxcolumn}[1]{>{\RaggedRight\arraybackslash}p{#1}}
\setlength{\parindent}{0pt}\setlength{\parskip}{5pt}\setlength{\emergencystretch}{2em}
\setlength{\tabcolsep}{4pt}\renewcommand{\arraystretch}{1.12}
\setlist{nosep,leftmargin=1.4em,topsep=3pt}
\titleformat{\section}{\large\bfseries\color{ink}}{\thesection}{.6em}{}
\titleformat{\subsection}{\normalsize\bfseries\color{ink}}{\thesubsection}{.6em}{}
\titlespacing*{\section}{0pt}{8pt}{6pt}\titlespacing*{\subsection}{0pt}{7pt}{4pt}
\pagestyle{fancy}\fancyhf{}
\fancyhead[L]{\small\color{gray}GEOI-BQ: BỔ SUNG BENCHMARK S4--S8}
\fancyhead[R]{\small\color{gray}09/11/2026}\fancyfoot[C]{\small\thepage}
\renewcommand{\headrulewidth}{.3pt}
\renewcommand{\contentsname}{Mục lục}\setcounter{tocdepth}{2}
\hypersetup{pdftitle={GeoI-BQ: bổ sung benchmark S4--S8},pdfauthor={}}
\newcommand{\PrimaryBA}{47{,}50}
\newcommand{\PrimaryFuture}{50{,}83}

\begin{document}
\pagenumbering{roman}
{\setlength{\parskip}{0pt}\tableofcontents}
\clearpage\pagenumbering{arabic}
\begin{center}
{\LARGE\bfseries\color{ink}Bổ sung đánh giá cho các scenario S4--S8}\\[3pt]
{\large GeoI-BQ v1 và năm phương án đối chứng}\\[2pt]
{\small Kết quả theo tác vụ; giữ riêng proxy, kiểm tra điều kiện và mục tiêu chưa đo được}
\end{center}

\section{Phạm vi và kết quả liên kết phiên S4}
\subsection{Protocol chung}
Sáu phương pháp: Raw GPS, REM-Z K1, TransProtect*, Semantic correlation*, Fake-query insertion*
và GeoI-BQ v1. Dấu * là adaptation đã audit của paper. Giữ nguyên output, GPS, Q và cấu hình đã đóng băng;
không lấy mẫu lại để cải thiện điểm. Dùng cùng 48 family: 20 TRAIN / 8 SELECTION / 20 TEST,
hai chuyến khác nhau/family, ba draws tại TEST, tương ứng 120 trip-draw/arm.

Fit attacker trên TRAIN, chọn theo từng arm/task/loss trên SELECTION, sau đó đánh giá TEST.
CI95\% dùng 10.000 bootstrap theo nhóm phụ thuộc. S5/S6 dùng family; S4/S8 dùng block gồm hai family
vì cặp âm hoặc partner không liên quan dùng hai family. Mỗi TEST có 10 block. Toàn bộ 66 tác vụ vòng đầu
và 24 tác vụ S8 vòng fusion hoàn tất. Lần kiểm tra gate quá chặt được giữ ở phụ lục D.
Đây là mở rộng trên cohort đã xem.

S1--S3/S9--S10 và Recall 99,595\% trong báo cáo 28/10 giữ nguyên nguồn. Bản này bổ sung tác vụ;
không coi utility toàn phiên cũ là một phép đo mới ở mọi cutoff. Formal privacy vẫn có các giả thiết
về common support, clock, bit đều và ledger; sampler Geo-I số thực/full restore chưa được chứng nhận đầy đủ.

\subsection{S4: đo liên kết phiên tới cùng nhóm tuyến}
\textbf{Target đo được:} hai phiên ngày 7/8 có thuộc cùng route family hay không. Ghép các family thành
block không giao nhau; mỗi block/draw có hai cặp cùng family và hai cặp chéo. TEST có 120 cặp cân bằng,
không dùng hai draws của cùng một chuyến làm hai phiên dương độc lập.

Attacker dùng \textbf{ExtraTrees96, logistic, kNN-1/5}, khoảng cách giữa các tập điểm ở đầu/cuối và uniform;
ngưỡng được chọn trên SELECTION. Feature đối xứng, chỉ đọc tọa độ/thời gian công khai. BA thấp hơn có lợi
cho unlinkability trong tác vụ này; mốc cân bằng là 50\%.
\begin{small}\par\smallskip\noindent\begin{tabularx}{\linewidth}{@{}P{3cm}XXX@{}}\toprule
\textbf{Phương pháp} & \textbf{BA\% [CI95\%]} & \textbf{AUC theo block [CI95\%]} & \textbf{Macro-F1} \\
\midrule
Raw GPS & 100{,}00 [100{,}00; 100{,}00] & 1{,}000 [1{,}000; 1{,}000] & 1{,}000 \\
REM-Z K1 & 91{,}67 [80{,}83; 98{,}33] & 0{,}956 [0{,}867; 1{,}000] & 0{,}916 \\
TransProtect* & 99{,}17 [97{,}50; 100{,}00] & 1{,}000 [1{,}000; 1{,}000] & 0{,}992 \\
Semantic correlation* & 100{,}00 [100{,}00; 100{,}00] & 1{,}000 [1{,}000; 1{,}000] & 1{,}000 \\
Fake-query insertion* & 93{,}33 [83{,}33; 100{,}00] & 1{,}000 [1{,}000; 1{,}000] & 0{,}933 \\
\textbf{GeoI-BQ v1} & 47{,}50 [40{,}00; 54{,}17] & 0{,}450 [0{,}386; 0{,}514] & 0{,}473 \\
\bottomrule\end{tabularx}\par\smallskip
\end{small}

GeoI-BQ có BA \PrimaryBA\%, gần mức đoán cân bằng; AUC dưới 0,5 không được diễn giải là bảo vệ
``tốt hơn ngẫu nhiên''. Các controls chứa tín hiệu liên kết mạnh hơn trong bank hiện tại.
AUC trong bảng là trung bình theo block với CI; AUC pooled và toàn bộ probabilities được giữ trong artifact.

\textbf{Giới hạn target.} Cohort chưa có ground truth người và phương tiện vật lý độc lập với route family.
Vì vậy, person re-identification và vehicle re-identification là \textbf{N/A}; không đổi tên điểm liên kết tuyến
thành kết quả bảo vệ danh tính người/xe. Account, biển số hoặc thiết bị trực tiếp cũng không được mô hình hóa.

\clearpage
\section{S5 và S6: dự đoán tương lai từ prefix}
\subsection{S5: cạnh ngoài kế tiếp tại fork hai lựa chọn}
Hai cạnh/đích ứng viên là auxiliary context công khai, chọn trước private route choice. Tên cạnh chỉ dùng
để xuất dự đoán; feature dùng hình dạng, khoảng cách và hướng. Target cạnh phải xuất hiện \emph{sau} cutoff
trong native SUMO. Attacker chỉ nhận các packet đã gửi tới cutoff, gồm cả fake packets.

Hai clock công khai được giữ: trước fork và cửa sổ đang rẽ. Không bổ sung GPS ở giữa các mốc gửi 20 s.
Bank gồm \textbf{candidate ExtraTrees96, kNN-1/5, geometry/motion và uniform}. Accuracy càng thấp thì
attacker càng khó chọn đúng cạnh trong bài toán hai lựa chọn này.
\begin{small}\par\smallskip\noindent\begin{tabularx}{\linewidth}{@{}P{3.2cm}XX@{}}\toprule
\textbf{Phương pháp} & \textbf{Trước fork: Acc\% [CI95\%]} & \textbf{Cửa sổ rẽ: Acc\% [CI95\%]} \\
\midrule
Raw GPS & 50{,}00 [50{,}00; 50{,}00] & 50{,}00 [50{,}00; 50{,}00] \\
REM-Z K1 & 50{,}00 [50{,}00; 50{,}00] & 50{,}00 [50{,}00; 50{,}00] \\
TransProtect* & 50{,}00 [50{,}00; 50{,}00] & 50{,}00 [50{,}00; 50{,}00] \\
Semantic correlation* & 50{,}00 [50{,}00; 50{,}00] & 50{,}00 [50{,}00; 50{,}00] \\
Fake-query insertion* & 44{,}17 [38{,}33; 50{,}00] & 50{,}00 [43{,}33; 57{,}50] \\
\textbf{GeoI-BQ v1} & 52{,}50 [48{,}33; 57{,}50] & 50{,}83 [45{,}83; 55{,}83] \\
\bottomrule\end{tabularx}\par\smallskip
\end{small}

\textbf{Phép đo S5 chưa có positive control mạnh.} Raw đạt 50\% ở cả hai pha; tại cửa sổ đang rẽ,
packet Raw cuối đã quan sát không nằm trên lane rẽ cần phân biệt trong các family TEST.
Lịch gửi 20 s đã bỏ lỡ khoảng rẽ ngắn. Vì vậy, điểm gần 50\% phần lớn phản ánh thông tin chưa được quan sát,
chưa chứng minh một cải thiện privacy của mô hình. Giữ đủ family và kết quả này, không lọc các trường hợp
để làm Raw tốt hơn hoặc proposed đẹp hơn. Mục tiêu next-route không giới hạn vẫn chưa được đánh giá.

\subsection{S6: đích đến sau khi đã có tín hiệu rẽ}
Thêm cutoff công khai $20\lceil t_{\rm turn}/20\rceil$, là mốc gửi đầu tiên từ clock đang rẽ trở đi.
Đích đến thật vẫn ở tương lai. Mỗi family có hai chuyến chọn hai đích cân bằng; baseline uniform là 50\%.
Không cung cấp route tương lai, routine label hoặc GPS tương lai cho attacker.
\begin{small}\par\smallskip\noindent\begin{tabularx}{\linewidth}{@{}P{3.1cm}XXX@{}}\toprule
\textbf{Phương pháp} & \textbf{Acc\% [CI95\%]} & \textbf{Hit100\% [CI95\%]} & \textbf{MAE m [CI95\%]} \\
\midrule
Raw GPS & 100{,}00 [100{,}00; 100{,}00] & 100{,}00 [100{,}00; 100{,}00] & 2{,}4 [2{,}4; 2{,}4] \\
REM-Z K1 & 50{,}00 [50{,}00; 50{,}00] & 50{,}00 [50{,}00; 50{,}00] & 572{,}7 [497{,}3; 647{,}8] \\
TransProtect* & 95{,}00 [90{,}83; 98{,}33] & 95{,}00 [90{,}83; 98{,}33] & 63{,}5 [17{,}7; 119{,}5] \\
Semantic correlation* & 95{,}00 [87{,}50; 100{,}00] & 95{,}00 [87{,}50; 100{,}00] & 63{,}9 [2{,}4; 156{,}2] \\
Fake-query insertion* & 99{,}17 [97{,}50; 100{,}00] & 99{,}17 [97{,}50; 100{,}00] & 11{,}3 [2{,}4; 29{,}1] \\
\textbf{GeoI-BQ v1} & 50{,}83 [44{,}17; 57{,}50] & 50{,}83 [44{,}17; 57{,}50] & 568{,}9 [458{,}8; 677{,}0] \\
\bottomrule\end{tabularx}\par\smallskip
\end{small}

Raw đạt 100\% và các adaptation phần lớn suy ra đúng đích; GeoI-BQ khoảng \PrimaryFuture\%,
REM-Z ở mức 50\%. Phép thử này có positive control phân biệt được target, nên có giá trị hơn pha trước fork.
Tuy nhiên, đây là chọn giữa \emph{hai đích đã biết}, không phải dự đoán một đích bất kỳ trong thành phố.

Hit100 trùng accuracy chọn đích trong workload này vì hai đích cách xa và đích đúng chỉ lệch ít mét so
với điểm dừng native. MAE được tính từ đích ứng viên được chọn tới GPS cuối. Brier/log loss đo chất lượng
probability, không gọi là ASR của paper khác. Các pha và CIs đầy đủ ở phụ lục A/B.
S6 dùng sáu chuyến lịch sử có bảo vệ chưa được chạy: outputs hiện tại chỉ bao gồm ngày 7/8.

\clearpage
\section{S7 và S8: nội dung truy vấn và kênh tương quan}
\subsection{S7: khả năng phân biệt nội dung khi giữ cùng trajectory}
Giữ cùng GPS/Q, thời điểm gửi và phản hồi provider; thay bốn purposes $\times$ sáu categories ở local.
Kiểm tra giao diện client đang dùng với một provider thử nghiệm trả phản hồi rỗng cố định;
requests vẫn hỏi mọi category ở L60. Reply-mask signatures của output thực được xác thực riêng.
Đây là kiểm tra giao diện/payload, không phải chạy lại POI utility hay suy luận ý định tự nhiên.
\begin{small}\par\smallskip\noindent\begin{tabularx}{\linewidth}{@{}P{3.2cm}XXXX@{}}\toprule
\textbf{Phương pháp} & \textbf{Wire bằng nhau} & \textbf{Purpose\%} & \textbf{Category\%} & \textbf{Joint\%} \\
\midrule
Raw GPS & 100{,}00 & 25{,}00 & 16{,}67 & 4{,}17 \\
REM-Z K1 & 100{,}00 & 25{,}00 & 16{,}67 & 4{,}17 \\
TransProtect* & 100{,}00 & 25{,}00 & 16{,}67 & 4{,}17 \\
Semantic correlation* & 100{,}00 & 25{,}00 & 16{,}67 & 4{,}17 \\
Fake-query insertion* & 100{,}00 & 25{,}00 & 16{,}67 & 4{,}17 \\
\textbf{GeoI-BQ v1} & 100{,}00 & 25{,}00 & 16{,}67 & 4{,}17 \\
\bottomrule\end{tabularx}\par\smallskip
\end{small}

Wire bằng nhau ở tất cả 1.056 streams, tương ứng 25.344 biến thể điều kiện. Với phân bố counterfactual
cân bằng, độ đoán đúng tối ưu trong một lớp wire bằng nhau là purpose $1/4$, category $1/6$, joint $1/24$.
Control thêm purpose/category trực tiếp vào request khôi phục 100\% nhãn joint.
Đây là \emph{conditional Bayes vulnerability}, không phải accuracy của attacker học trên hành vi người thật.

Cả sáu arm dùng chung cover wrapper, nên S7 là tính chất của giao diện chung, không là ưu thế riêng trước
implementation gốc của các paper. Ý định tương quan tuyến, tài khoản hoặc lúc mở phiên vẫn chưa được đo.

\subsection{S8: suy luận vị trí từ hai kênh có shared prefix}
Dùng hai phiên ngày 7/8 khác nhau, căn theo clock tương đối, cắt trước fork tại một tick công khai.
Native GPS của hai phiên giống nhau trên prefix này. \textbf{Đây là mô hình tương quan tổng hợp sau căn clock},
không là hai xe đồng hành chạy đồng thời trong SUMO. Attacker được cấp liên kết giữa hai kênh.

Mỗi ô là \textbf{Hit100\% / MAE m}. Target: chỉ thấy kênh mục tiêu; +protected: thêm kênh đã bảo vệ cùng
shared prefix; +raw: cấp công khai GPS của kênh tương quan; +unrelated: partner bảo vệ từ family khác.
Bảng dùng bank fusion V2 được công bố ở Phần 4.
Hit100 và MAE được chọn attacker riêng theo từng loss, nên hai số trong một ô có thể từ hai estimators khác nhau.
\begin{small}\par\smallskip\noindent\begin{tabularx}{\linewidth}{@{}P{3.2cm}XXXX@{}}\toprule
\textbf{Phương pháp} & \textbf{Target} & \textbf{+ protected} & \textbf{+ raw} & \textbf{+ unrelated} \\
\midrule
Raw GPS & 100{,}00 / 0{,}0 & 100{,}00 / 0{,}0 & 100{,}00 / 0{,}0 & 100{,}00 / 0{,}0 \\
REM-Z K1 & 10{,}83 / 359{,}7 & 6{,}67 / 300{,}8 & 100{,}00 / 0{,}0 & 10{,}83 / 359{,}7 \\
TransProtect* & 60{,}83 / 102{,}8 & 64{,}17 / 96{,}8 & 98{,}33 / 0{,}0 & 60{,}83 / 102{,}8 \\
Semantic correlation* & 45{,}00 / 101{,}4 & 75{,}00 / 46{,}2 & 100{,}00 / 0{,}0 & 45{,}00 / 101{,}4 \\
Fake-query insertion* & 16{,}67 / 1857{,}4 & 68{,}33 / 1857{,}4 & 100{,}00 / 0{,}0 & 16{,}67 / 1857{,}4 \\
\textbf{GeoI-BQ v1} & 0{,}00 / 2246{,}7 & 0{,}00 / 2226{,}8 & 100{,}00 / 0{,}0 & 0{,}00 / 2246{,}7 \\
\bottomrule\end{tabularx}\par\smallskip
\end{small}

GeoI-BQ không có hit trong mẫu ở target/+protected, nhưng không suy từ đó rủi ro bằng 0.
Khi partner Raw được công khai, estimator đọc trực tiếp partner định vị được 100\% target; attacker đã chọn
trên SELECTION đạt 98,33--100\% ở TEST. GPS tương quan đủ để tiết lộ target trong mô hình này.
Điều đó thể hiện giới hạn của việc bảo vệ một người khi có side information rất mạnh.
Một cận single-session không tự trở thành chứng nhận cho joint transcript của nhiều người/phiên.

\clearpage
\section{Attacker mạnh hơn, kết luận và phạm vi bằng chứng}
\subsection{S8 V1 sang V2: giữ lại quá trình sửa bank attacker}
Sau khi đọc V1, audit nhận ra bank thiếu phép lấy trung bình hai kênh và tìm cặp Q gần nhau nhất.
V2 giữ nguyên toàn bộ V1 và thêm centroid/median/velocity mixtures với trọng số $1/4,1/2,3/4$;
midpoint cặp gần nhất và trung bình hai/ba cặp gần nhất. Các phép này chỉ dùng hai prefix công khai.
Nếu mỗi tập đều chứa cùng GPS thật, việc đối chiếu hai tập có thể làm lộ thành viên đó.

V2 được chốt trước khi chấm điểm riêng, chọn lại attacker trên SELECTION; đây là \textbf{post-V1 diagnostic
trên dữ liệu đã xem}, không phải confirmation mới. Bảng giữ cả hai vòng, view +protected:
\begin{small}\par\smallskip\noindent\begin{tabularx}{\linewidth}{@{}P{3.1cm}XXXX@{}}\toprule
\textbf{Phương pháp} & \textbf{V1 Hit\%} & \textbf{V2 Hit\%} & \textbf{V1 MAE m} & \textbf{V2 MAE m} \\
\midrule
Raw GPS & 100{,}00 & 100{,}00 & 0{,}0 & 0{,}0 \\
REM-Z K1 & 10{,}83 & 6{,}67 & 359{,}7 & 300{,}8 \\
TransProtect* & 60{,}83 & 64{,}17 & 102{,}8 & 96{,}8 \\
Semantic correlation* & 45{,}00 & 75{,}00 & 101{,}4 & 46{,}2 \\
Fake-query insertion* & 16{,}67 & 68{,}33 & 1857{,}4 & 1857{,}4 \\
\textbf{GeoI-BQ v1} & 0{,}00 & 0{,}00 & 2226{,}8 & 2226{,}8 \\
\bottomrule\end{tabularx}\par\smallskip
\end{small}

Fusion làm attacker định vị tốt hơn rõ ở Semantic*/Fake-query* trong readout này. Nó không bảo đảm
mọi metric TEST đều tăng: selector được giữ theo SELECTION và chọn riêng cho Hit/MAE.
Kết quả không được chọn bằng cách lấy attacker có score đẹp nhất sau khi xem TEST.

\subsection{Những attacker thực sự được chọn cho GeoI-BQ}
\begin{small}\par\smallskip\noindent\begin{tabularx}{\linewidth}{@{}P{3cm}X@{}}\toprule
\textbf{Tác vụ GeoI-BQ} & \textbf{Attacker đã chọn trên SELECTION} \\
\midrule
S4 association & extra\_trees96\, @\, 0.5 \\
S5 turn & curve\_min\_0.02 \\
S6 post-turn & curve\_min\_0.02 \\
S8 Target & Hit: centroid; MAE: learned\_knn5 \\
S8 + protected & Hit: centroid; MAE: learned\_knn5 \\
S8 + raw & Hit: partner\_centroid; MAE: partner\_centroid \\
S8 + unrelated & Hit: centroid; MAE: target\_only\_knn5 \\
\bottomrule\end{tabularx}\par\smallskip
\end{small}
\texttt{curve\_min} dùng khoảng cách từ Q tới hình dạng rẽ công khai; các temperature được chốt trong bank.
Một selected model là kết quả của finite bank, không chứng minh attacker tối ưu đã thất bại.

\subsection{Kết luận có thể trình bày}
\begin{enumerate}
\item Đã bổ sung phép đo tái lập cho liên kết phiên S4, fork S5, đích S6, payload S7 và kênh tương quan S8,
với target/metrics riêng thay vì điền cùng Hit100 cho mọi scenario.
\item Bằng chứng S4 là proxy liên kết tuyến; S6 sau rẽ có positive control mạnh. S5 hiện có hạn chế về
khả năng quan sát do cadence, nên chưa là bằng chứng superiority.
\item S7 xác nhận nhu cầu local không thay payload trong điều kiện cố định. S8 cho thấy raw side information
có thể phá bảo vệ vị trí khi hai trajectory tương quan hoàn hảo; kết quả +protected cần đọc cùng scope tổng hợp.
\item Chứng minh privacy toàn transcript tiếp tục là nền tảng. Benchmark đo hiệu quả trong những kênh và
workload cụ thể, không thay các giả thiết của định lý hoặc chứng nhận triển khai.
\end{enumerate}
Các mục còn N/A được giữ công khai: identity người/xe thật; S6 protected history sáu chuyến;
S7 natural intent; S8 co-travel/companion detection thực tế. Đây là giới hạn dữ liệu/tác vụ, không phải score 0.

\clearpage
\section*{Phụ lục A. S5: xác suất dự đoán và CI đầy đủ}
\addcontentsline{toc}{section}{Phụ lục A. S5: xác suất dự đoán và CI đầy đủ}
Accuracy là chọn đúng ID cạnh trong hai ứng viên công khai. CI95\% theo 20 TEST family, mỗi family
hai chuyến $\times$ ba draws. Brier là $\mathbb E(p_1-y)^2$; log loss dùng log tự nhiên.
\begin{small}\par\smallskip\noindent\begin{tabularx}{\linewidth}{@{}P{3cm}P{2.1cm}X P{1cm}P{1cm}@{}}\toprule
\textbf{Phương pháp} & \textbf{Pha} & \textbf{Acc\% [CI95\%]} & \textbf{Brier} & \textbf{Log loss} \\
\midrule
Raw GPS & Trước fork & 50{,}00 [50{,}00; 50{,}00] & 0{,}250 & 0{,}693 \\
Raw GPS & Cửa sổ rẽ & 50{,}00 [50{,}00; 50{,}00] & 0{,}250 & 0{,}693 \\
REM-Z K1 & Trước fork & 50{,}00 [50{,}00; 50{,}00] & 0{,}250 & 0{,}693 \\
REM-Z K1 & Cửa sổ rẽ & 50{,}00 [50{,}00; 50{,}00] & 0{,}250 & 0{,}693 \\
TransProtect* & Trước fork & 50{,}00 [50{,}00; 50{,}00] & 0{,}250 & 0{,}693 \\
TransProtect* & Cửa sổ rẽ & 50{,}00 [50{,}00; 50{,}00] & 0{,}250 & 0{,}693 \\
Semantic correlation* & Trước fork & 50{,}00 [50{,}00; 50{,}00] & 0{,}250 & 0{,}693 \\
Semantic correlation* & Cửa sổ rẽ & 50{,}00 [50{,}00; 50{,}00] & 0{,}252 & 0{,}697 \\
Fake-query insertion* & Trước fork & 44{,}17 [38{,}33; 50{,}00] & 0{,}315 & 0{,}853 \\
Fake-query insertion* & Cửa sổ rẽ & 50{,}00 [43{,}33; 57{,}50] & 0{,}373 & 5{,}372 \\
\textbf{GeoI-BQ v1} & Trước fork & 52{,}50 [48{,}33; 57{,}50] & 0{,}259 & 0{,}712 \\
\textbf{GeoI-BQ v1} & Cửa sổ rẽ & 50{,}83 [45{,}83; 55{,}83] & 0{,}259 & 0{,}714 \\
\bottomrule\end{tabularx}\par\smallskip
\end{small}
Raw trước fork có cùng prefix cho hai lựa chọn nên accuracy 50\% là mốc thiếu thông tin nội tại.
Tại cửa sổ rẽ, clock native đúng vẫn chưa bảo đảm có packet đã phát trong cửa sổ đó. Các row và số lượng
tọa độ/fake packets quan sát được được lưu theo từng sample; không dùng last-N packet thay cho cutoff.

\section*{Nguyên tắc đọc những điểm gần ngẫu nhiên}
\begin{itemize}
\item Phải có Raw/positive control đủ tín hiệu trước khi ghi nhận đóng góp của defense.
\item Điểm dưới 50\% không tự là privacy tốt hơn; có thể là mất calibration hoặc biến thiên mẫu.
\item CI chứa mức ngẫu nhiên chỉ mô tả readout của bank này; nó không chứng nhận mọi attacker đều vô lợi.
\end{itemize}

\clearpage
\section*{Phụ lục B. S6: tất cả pha, xác suất và sai số đích}
\addcontentsline{toc}{section}{Phụ lục B. S6: tất cả pha, xác suất và sai số đích}
Phân biệt ba phạm vi quan sát; không dùng cả chuyến đã đi để dự đoán đích vẫn gọi là ``tương lai''.
Eligibility kiểm tra arrival native nằm sau cutoff ở tất cả rows.
\begin{footnotesize}\par\smallskip\noindent\begin{tabularx}{\linewidth}{@{}P{3cm}P{2.1cm}X P{1cm}P{1cm}X@{}}\toprule
\textbf{Phương pháp} & \textbf{Pha} & \textbf{Acc\% [CI95\%]} & \textbf{Brier} & \textbf{Log loss} & \textbf{MAE m [CI95\%]} \\
\midrule
Raw GPS & Trước fork & 50{,}00 [50{,}00; 50{,}00] & 0{,}250 & 0{,}693 & 572{,}7 [497{,}3; 647{,}8] \\
Raw GPS & Cửa sổ rẽ & 50{,}00 [50{,}00; 50{,}00] & 0{,}250 & 0{,}693 & 572{,}7 [497{,}3; 647{,}8] \\
Raw GPS & Sau rẽ & 100{,}00 [100{,}00; 100{,}00] & 0{,}034 & 0{,}181 & 2{,}4 [2{,}4; 2{,}4] \\
REM-Z K1 & Trước fork & 50{,}00 [50{,}00; 50{,}00] & 0{,}250 & 0{,}693 & 572{,}7 [497{,}3; 647{,}8] \\
REM-Z K1 & Cửa sổ rẽ & 50{,}00 [50{,}00; 50{,}00] & 0{,}250 & 0{,}693 & 572{,}7 [497{,}3; 647{,}8] \\
REM-Z K1 & Sau rẽ & 50{,}00 [50{,}00; 50{,}00] & 0{,}250 & 0{,}693 & 572{,}7 [497{,}3; 647{,}8] \\
TransProtect* & Trước fork & 50{,}00 [50{,}00; 50{,}00] & 0{,}250 & 0{,}693 & 572{,}7 [497{,}3; 647{,}8] \\
TransProtect* & Cửa sổ rẽ & 50{,}00 [50{,}00; 50{,}00] & 0{,}250 & 0{,}693 & 572{,}7 [497{,}3; 647{,}8] \\
TransProtect* & Sau rẽ & 95{,}00 [90{,}83; 98{,}33] & 0{,}168 & 0{,}518 & 63{,}5 [17{,}7; 119{,}5] \\
Semantic correlation* & Trước fork & 50{,}00 [50{,}00; 50{,}00] & 0{,}250 & 0{,}693 & 572{,}7 [497{,}3; 647{,}8] \\
Semantic correlation* & Cửa sổ rẽ & 50{,}00 [50{,}00; 50{,}00] & 0{,}252 & 0{,}697 & 572{,}7 [497{,}3; 647{,}8] \\
Semantic correlation* & Sau rẽ & 95{,}00 [87{,}50; 100{,}00] & 0{,}128 & 0{,}436 & 63{,}9 [2{,}4; 156{,}2] \\
Fake-query insertion* & Trước fork & 44{,}17 [38{,}33; 50{,}00] & 0{,}315 & 0{,}853 & 649{,}4 [528{,}1; 784{,}3] \\
Fake-query insertion* & Cửa sổ rẽ & 50{,}00 [43{,}33; 57{,}50] & 0{,}373 & 5{,}372 & 554{,}0 [461{,}6; 643{,}4] \\
Fake-query insertion* & Sau rẽ & 99{,}17 [97{,}50; 100{,}00] & 0{,}119 & 0{,}416 & 11{,}3 [2{,}4; 29{,}1] \\
\textbf{GeoI-BQ v1} & Trước fork & 52{,}50 [48{,}33; 57{,}50] & 0{,}259 & 0{,}712 & 540{,}5 [452{,}8; 627{,}4] \\
\textbf{GeoI-BQ v1} & Cửa sổ rẽ & 50{,}83 [45{,}83; 55{,}83] & 0{,}259 & 0{,}714 & 558{,}5 [469{,}3; 646{,}7] \\
\textbf{GeoI-BQ v1} & Sau rẽ & 50{,}83 [44{,}17; 57{,}50] & 0{,}259 & 0{,}712 & 568{,}9 [458{,}8; 677{,}0] \\
\bottomrule\end{tabularx}\par\smallskip
\end{footnotesize}
Candidate accuracy và Hit100 bằng nhau trong trial hai đích này; Brier/log loss vẫn có thể khác giữa
các phương pháp có cùng accuracy. Không gọi số MAE này là kết quả của một attacker đích đến không giới hạn.

\clearpage
\section*{Phụ lục C. S8 V2: CI và chênh lệch ghép theo block}
\addcontentsline{toc}{section}{Phụ lục C. S8 V2: CI và chênh lệch ghép theo block}
Đơn vị ΔHit là điểm phần trăm; ΔMAE theo mét. Chênh lệch là view trừ target-only, cùng session/draw.
CI95\% ghép theo 10 block để giữ phụ thuộc khi hai family trao đổi unrelated partners.
\begin{footnotesize}\par\smallskip\noindent\begin{tabularx}{\linewidth}{@{}P{2.8cm}P{1.5cm}XXXX@{}}\toprule
\textbf{Phương pháp} & \textbf{View} & \textbf{Hit\% [CI]} & \textbf{MAE m [CI]} & \textbf{ΔHit pp [CI]} & \textbf{ΔMAE m [CI]} \\
\midrule
Raw GPS & Target & 100{,}00\newline [100{,}00; 100{,}00] & 0{,}0\newline [0{,}0; 0{,}0] & 0{,}00\newline [0{,}00; 0{,}00] & 0{,}0\newline [0{,}0; 0{,}0] \\
Raw GPS & + protected & 100{,}00\newline [100{,}00; 100{,}00] & 0{,}0\newline [0{,}0; 0{,}0] & 0{,}00\newline [0{,}00; 0{,}00] & 0{,}0\newline [0{,}0; 0{,}0] \\
Raw GPS & + raw & 100{,}00\newline [100{,}00; 100{,}00] & 0{,}0\newline [0{,}0; 0{,}0] & 0{,}00\newline [0{,}00; 0{,}00] & 0{,}0\newline [0{,}0; 0{,}0] \\
Raw GPS & + unrelated & 100{,}00\newline [100{,}00; 100{,}00] & 0{,}0\newline [0{,}0; 0{,}0] & 0{,}00\newline [0{,}00; 0{,}00] & 0{,}0\newline [0{,}0; 0{,}0] \\
REM-Z K1 & Target & 10{,}83\newline [5{,}00; 16{,}67] & 359{,}7\newline [297{,}6; 424{,}9] & 0{,}00\newline [0{,}00; 0{,}00] & 0{,}0\newline [0{,}0; 0{,}0] \\
REM-Z K1 & + protected & 6{,}67\newline [0{,}00; 15{,}00] & 300{,}8\newline [247{,}9; 357{,}9] & -4{,}17\newline [-11{,}67; 2{,}50] & -58{,}9\newline [-74{,}7; -43{,}8] \\
REM-Z K1 & + raw & 100{,}00\newline [100{,}00; 100{,}00] & 0{,}0\newline [0{,}0; 0{,}0] & 89{,}17\newline [83{,}33; 95{,}00] & -359{,}7\newline [-424{,}9; -297{,}6] \\
REM-Z K1 & + unrelated & 10{,}83\newline [5{,}00; 16{,}67] & 359{,}7\newline [297{,}6; 424{,}9] & 0{,}00\newline [0{,}00; 0{,}00] & 0{,}0\newline [0{,}0; 0{,}0] \\
TransProtect* & Target & 60{,}83\newline [51{,}67; 70{,}00] & 102{,}8\newline [84{,}8; 121{,}9] & 0{,}00\newline [0{,}00; 0{,}00] & 0{,}0\newline [0{,}0; 0{,}0] \\
TransProtect* & + protected & 64{,}17\newline [56{,}67; 71{,}67] & 96{,}8\newline [78{,}8; 116{,}2] & 3{,}33\newline [0{,}00; 6{,}67] & -5{,}9\newline [-8{,}5; -3{,}6] \\
TransProtect* & + raw & 98{,}33\newline [95{,}83; 100{,}00] & 0{,}0\newline [0{,}0; 0{,}0] & 37{,}50\newline [27{,}50; 47{,}50] & -102{,}8\newline [-121{,}9; -84{,}8] \\
TransProtect* & + unrelated & 60{,}83\newline [51{,}67; 70{,}00] & 102{,}8\newline [84{,}8; 121{,}9] & 0{,}00\newline [0{,}00; 0{,}00] & 0{,}0\newline [0{,}0; 0{,}0] \\
Semantic correlation* & Target & 45{,}00\newline [20{,}00; 70{,}00] & 101{,}4\newline [76{,}7; 132{,}1] & 0{,}00\newline [0{,}00; 0{,}00] & 0{,}0\newline [0{,}0; 0{,}0] \\
Semantic correlation* & + protected & 75{,}00\newline [60{,}00; 90{,}00] & 46{,}2\newline [29{,}3; 62{,}5] & 30{,}00\newline [10{,}00; 50{,}00] & -55{,}3\newline [-95{,}8; -23{,}9] \\
Semantic correlation* & + raw & 100{,}00\newline [100{,}00; 100{,}00] & 0{,}0\newline [0{,}0; 0{,}0] & 55{,}00\newline [30{,}00; 80{,}00] & -101{,}4\newline [-132{,}1; -76{,}7] \\
Semantic correlation* & + unrelated & 45{,}00\newline [20{,}00; 70{,}00] & 101{,}4\newline [76{,}7; 132{,}1] & 0{,}00\newline [0{,}00; 0{,}00] & 0{,}0\newline [0{,}0; 0{,}0] \\
Fake-query insertion* & Target & 16{,}67\newline [5{,}00; 31{,}67] & 1857{,}4\newline [1526{,}6; 2225{,}8] & 0{,}00\newline [0{,}00; 0{,}00] & 0{,}0\newline [0{,}0; 0{,}0] \\
Fake-query insertion* & + protected & 68{,}33\newline [56{,}67; 80{,}00] & 1857{,}4\newline [1526{,}6; 2225{,}8] & 51{,}67\newline [46{,}67; 56{,}67] & 0{,}0\newline [0{,}0; 0{,}0] \\
Fake-query insertion* & + raw & 100{,}00\newline [100{,}00; 100{,}00] & 0{,}0\newline [0{,}0; 0{,}0] & 83{,}33\newline [68{,}33; 95{,}00] & -1857{,}4\newline [-2225{,}8; -1526{,}6] \\
Fake-query insertion* & + unrelated & 16{,}67\newline [5{,}00; 31{,}67] & 1857{,}4\newline [1526{,}6; 2225{,}8] & 0{,}00\newline [0{,}00; 0{,}00] & 0{,}0\newline [0{,}0; 0{,}0] \\
\textbf{GeoI-BQ v1} & Target & 0{,}00\newline [0{,}00; 0{,}00] & 2246{,}7\newline [2036{,}7; 2468{,}2] & 0{,}00\newline [0{,}00; 0{,}00] & 0{,}0\newline [0{,}0; 0{,}0] \\
\textbf{GeoI-BQ v1} & + protected & 0{,}00\newline [0{,}00; 0{,}00] & 2226{,}8\newline [1950{,}9; 2507{,}3] & 0{,}00\newline [0{,}00; 0{,}00] & -19{,}9\newline [-161{,}8; 111{,}7] \\
\textbf{GeoI-BQ v1} & + raw & 100{,}00\newline [100{,}00; 100{,}00] & 0{,}0\newline [0{,}0; 0{,}0] & 100{,}00\newline [100{,}00; 100{,}00] & -2246{,}7\newline [-2468{,}2; -2036{,}7] \\
\textbf{GeoI-BQ v1} & + unrelated & 0{,}00\newline [0{,}00; 0{,}00] & 2246{,}7\newline [2036{,}7; 2468{,}2] & 0{,}00\newline [0{,}00; 0{,}00] & 0{,}0\newline [0{,}0; 0{,}0] \\
\bottomrule\end{tabularx}\par\smallskip
\end{footnotesize}
CI [0;0] khi không có hit là giới hạn của percentile bootstrap trên mẫu, không là chứng minh xác suất hit bằng 0.
Positive control +raw là side information được cấp công khai, không phải GPS target bị đưa lén vào features.

\clearpage
\section*{Phụ lục D. Số quan sát, kiểm chứng và tái lập}
\addcontentsline{toc}{section}{Phụ lục D. Số quan sát, kiểm chứng và tái lập}
\textbf{Số tọa độ công khai trung bình mà attacker quan sát tại cutoff S8.}
Joint views tính cả hai kênh và mọi fake packets tới cutoff. Đây là số quan sát của attacker,
không phải thời gian mạng hoặc một lần benchmark utility/energy mới.
\begin{small}\par\smallskip\noindent\begin{tabularx}{\linewidth}{@{}P{3.2cm}XXXX@{}}\toprule
\textbf{Phương pháp} & \textbf{Target} & \textbf{+ protected} & \textbf{+ raw} & \textbf{+ unrelated} \\
\midrule
Raw GPS & 6{,}2 & 12{,}4 & 12{,}4 & 12{,}4 \\
REM-Z K1 & 6{,}2 & 12{,}4 & 12{,}4 & 12{,}4 \\
TransProtect* & 6{,}2 & 12{,}4 & 12{,}4 & 12{,}4 \\
Semantic correlation* & 31{,}0 & 62{,}0 & 37{,}2 & 62{,}0 \\
Fake-query insertion* & 36{,}8 & 73{,}5 & 43{,}0 & 73{,}5 \\
\textbf{GeoI-BQ v1} & 31{,}0 & 62{,}0 & 37{,}2 & 62{,}0 \\
\bottomrule\end{tabularx}\par\smallskip
\end{small}

Nguồn chính: \texttt{remaining\_scenarios\_20261109\_v1}; vòng sửa attacker:
\texttt{s8\_fusion\_20261109\_v2}, cùng trong \texttt{artifacts/benchmarks/}.
Nguồn cũ \texttt{paper\_comparison\_20261028\_v1} giữ nguyên, có source hash được xác thực.
Hai protocol ghi rõ trước/sau readout V1 và các giới hạn scope.

\textbf{Verifier.} Tái tính metrics/CIs từ predictions đã lưu, xác thực split/block, cutoff, cạnh/arrival tương lai,
nhãn riêng không vào feature API và ngân sách/source pins của mô hình. V2 lưu cả SELECTION predictions
để replay đúng selector, rồi mới so TEST. Tất cả failure, models và candidate predictions được giữ.

Từ repo root, dùng môi trường \texttt{requirements-paper-comparison.txt}:
\begin{quote}\small\ttfamily
python -m experiments.verify\_remaining\_scenarios\_20261109\\
python -m experiments.check\_s8\_fusion\_20261109
\end{quote}
Trong thư mục report:
\begin{quote}\small\ttfamily
python build\_report.py --check\\
tectonic --keep-logs report\_explained.tex\\
python validate\_report.py
\end{quote}
Sau khi render và xem toàn bộ PDF, validator với \texttt{--visual-reviewed} ghi receipt QA.
Các bảng được xuất từ JSON đã qua verifier; exporter không fit attacker, tạo lại Q hoặc sửa core Geo-I.
Lệnh \texttt{--write} của verifier chỉ dùng lần đầu để tạo receipt write-once.

\textbf{Gate kiểm tra đã sửa.} Verifier V2 ban đầu đòi model được chọn phải đạt Hit100=100\% khi có Raw
partner. Một fusion estimator cùng đạt 100\% trên SELECTION nhưng chỉ đạt 98,33\% trên TEST ở TransProtect*.
Đó là biến thiên của selector, không phải mất Raw side channel: estimator \texttt{partner\_centroid} vẫn đạt
100\%. Giữ failure đầu tiên và dùng checker độc lập để xác nhận cả available control lẫn selected score;
không sửa prediction, kết quả hoặc defense. Không gọi lại frozen verifier với gate sai để ghi đè receipt.

\textbf{Nguồn phương pháp.} Báo cáo 28/10 cung cấp kiến trúc/bốn mệnh đề privacy và utility;
hai protocol 09/11 cung cấp hợp đồng đánh giá bổ sung. Chứng minh không biến proxy identity hoặc
hai phiên căn clock thành ground truth người/xe hay co-travel thực tế. Để mở rộng những mục tiêu đó,
cần dữ liệu/quan sát tương ứng và phân tích nhiều kênh riêng.
\end{document}
